\documentclass[aspectratio=169,14pt]{beamer}

\usetheme{Madrid}
\usecolortheme{whale}

\definecolor{ctiblue}{RGB}{31,78,121}
\definecolor{ctiorange}{RGB}{196,89,17}
\definecolor{ctigreen}{RGB}{39,124,52}
\definecolor{ctired}{RGB}{180,30,30}
\definecolor{lightblue}{RGB}{230,240,250}

\setbeamercolor{structure}{fg=ctiblue}
\setbeamercolor{frametitle}{bg=ctiblue,fg=white}
\setbeamercolor{title}{bg=ctiblue,fg=white}
\setbeamercolor{block title}{bg=ctiblue,fg=white}
\setbeamercolor{block body}{bg=lightblue}
\setbeamercolor{block title alerted}{bg=ctired,fg=white}
\setbeamercolor{block title example}{bg=ctigreen,fg=white}

\setbeamertemplate{navigation symbols}{}
\setbeamertemplate{footline}[frame number]

\usepackage{amsmath,amssymb}
\usepackage{booktabs}
\usepackage{graphicx}

\DeclareMathOperator{\logit}{logit}

\title{The CTI Universal Law}
\subtitle{A First-Principles Visual Guide}
\author{Devansh}
\date{March 2026}

\begin{document}

% ===================== TITLE =====================
\begin{frame}
\titlepage
\end{frame}

% ===================== SLIDE 1 =====================
\begin{frame}{The Big Idea}
\begin{center}
\Large
There is a \textbf{single equation} governing\\
how well \textit{any} representation separates categories.\\[0.8cm]

\Huge\color{ctiblue}
$\logit(q) = \alpha\,\kappa - \beta\log(K{-}1) + C$\\[0.8cm]

\normalsize\color{black}
Derived from first principles.\\
Validated across NLP, vision, audio, and biological cortex.
\end{center}
\end{frame}

% ===================== SLIDE 2 =====================
\begin{frame}{What Are We Measuring?}
\begin{columns}[T]
\begin{column}{0.55\textwidth}
\textbf{The question}: How good are a model's internal representations at telling categories apart?

\vspace{0.3cm}
\textbf{The probe}: 1-Nearest Neighbor (1-NN)
\begin{itemize}
\item Find closest training point
\item Predict same label
\item No learnable parameters
\item Accuracy = pure geometry
\end{itemize}
\end{column}
\begin{column}{0.42\textwidth}
\begin{block}{Why not a linear probe?}
\small
Linear probes mix geometry with optimization. 1-NN removes this confounder: \textbf{accuracy IS geometry}.
\end{block}
\end{column}
\end{columns}
\end{frame}

% ===================== SLIDE 3 =====================
\begin{frame}{Key Quantity \#1: $q_{\mathrm{norm}}$}
\begin{block}{Chance-Normalized Accuracy}
\[
q_{\mathrm{norm}} = \frac{\mathrm{acc}_{1\mathrm{NN}} - 1/K}{1 - 1/K}
\]
\end{block}

\begin{itemize}
\item $q = 0$: at chance (no signal)
\item $q = 1$: perfect classification
\item Comparable across different $K$
\end{itemize}

\vspace{0.3cm}
\textbf{Example}: $K{=}10$, acc $= 55\%$ $\to$ $q = (0.55 - 0.10)/(0.90) = 0.50$
\end{frame}

% ===================== SLIDE 4 =====================
\begin{frame}{Key Quantity \#2: $\kappa_{\mathrm{nearest}}$}
\begin{block}{The Geometric Signal-to-Noise Ratio}
\[
\kappa_{\mathrm{nearest}}(k) = \min_{j \neq k} \frac{\|\boldsymbol{\mu}_j - \boldsymbol{\mu}_k\|}{\sigma_W \sqrt{d}}
\]
\end{block}

\begin{columns}[T]
\begin{column}{0.5\textwidth}
\textbf{Numerator:}\\
Distance to the \textbf{nearest} other class centroid (hardest competitor)

\vspace{0.3cm}
\textbf{Denominator:}\\
Expected within-class L2 radius (noise level)
\end{column}
\begin{column}{0.45\textwidth}
\begin{alertblock}{Why ``nearest''?}
\small
In a 1-NN race, you confuse with the \textbf{nearest} class. Far classes are irrelevant. This is the \textbf{bottleneck}.
\end{alertblock}
\end{column}
\end{columns}
\end{frame}

% ===================== SLIDE 5 =====================
\begin{frame}{$\kappa$ Intuition: Signal vs.\ Noise}
\begin{center}
\begin{tabular}{ccl}
\toprule
$\kappa$ & Regime & Meaning \\
\midrule
$\ll 1$ & Noise dominates & Classification at chance \\
$\sim 1$ & Signal $\approx$ noise & Transition zone \\
$\gg 1$ & Signal dominates & Near-perfect \\
\bottomrule
\end{tabular}
\end{center}

\vspace{0.5cm}
$\kappa$ is \textbf{dimensionless} --- comparable across models with different embedding dimensions.

\vspace{0.3cm}
Think of it like a ``contrast ratio'' for class separation.
\end{frame}

% ===================== SLIDE 6 =====================
\begin{frame}{The Gumbel Race: Step 1 --- 1-NN as a Race}
\begin{block}{The Competition}
\begin{align*}
M_+ &= \text{distance to nearest \textbf{same-class} point} \\
M_j &= \text{distance to nearest point from class } j
\end{align*}
\end{block}

\vspace{0.3cm}
\begin{center}
\Large
1-NN is correct $\iff$ $M_+ < \min_{j \neq k} M_j$
\end{center}

\vspace{0.3cm}
The correct class must \textbf{win the race} by having the smallest minimum distance.
\end{frame}

% ===================== SLIDE 7 =====================
\begin{frame}{The Gumbel Race: Step 2 --- Extreme Value Theory}
\textbf{Fisher--Tippett theorem (1928):} The minimum of $n$ i.i.d.\ random variables converges to one of three universal distributions:

\vspace{0.3cm}
\begin{enumerate}
\item \textbf{Gumbel} --- exponential tails (Gaussian, etc.) $\leftarrow$ \textcolor{ctiblue}{\textbf{Our case}}
\item Fr\'echet --- heavy tails
\item Weibull --- finite support
\end{enumerate}

\vspace{0.3cm}
\begin{block}{After EVT normalization}
\begin{align*}
M_+ &\sim \mathrm{Gumbel}(\mu_+, s) \\
M_j &\sim \mathrm{Gumbel}(\mu_j^{(-)}, s) \quad \text{shared scale } s
\end{align*}
\end{block}
\end{frame}

% ===================== SLIDE 8 =====================
\begin{frame}{The Gumbel Race: Step 3 --- The Magic Identity}
\begin{block}{The Gumbel Race Identity (EXACT)}
If $G_0, G_1, \ldots, G_{K-1}$ are independent Gumbels with common scale $s$:
\[
\boxed{P(G_0 < \min_{j \geq 1} G_j) = \frac{1}{1 + \sum_{j} \exp\!\bigl(-(\mu_j - \mu_0)/s\bigr)}}
\]
\end{block}

\vspace{0.5cm}
\begin{center}
\large
This is \textbf{exact}. Not an approximation.\\[0.3cm]
The probability is a \textbf{logistic function} of the location gaps.\\[0.3cm]
\textcolor{ctiblue}{\textbf{This is why the law has a sigmoid/logit form.}}
\end{center}
\end{frame}

% ===================== SLIDE 9 =====================
\begin{frame}{The Gumbel Race: Steps 4--5 --- Competition Structure}
\textbf{Apply to 1-NN:}
\[
q = \frac{1}{1 + \sum_{j \neq k} \exp(-\Delta_{k,j}/s)}
\]

\vspace{0.3cm}
\textbf{Two regimes:}

\begin{columns}[T]
\begin{column}{0.48\textwidth}
\begin{block}{Dense (ETF)}
All $K{-}1$ competitors equal:\\
$\sum = (K{-}1)e^{-\Delta/s}$\\
$\Rightarrow$ $\beta = 1$
\end{block}
\end{column}
\begin{column}{0.48\textwidth}
\begin{exampleblock}{Hierarchical (Real World)}
Only $\sqrt{K}$ active competitors:\\
$\sum = \sqrt{K{-}1}\,e^{-\Delta/s}$\\
$\Rightarrow$ $\beta = 0.5$
\end{exampleblock}
\end{column}
\end{columns}

\vspace{0.3cm}
\textbf{Empirical}: $\hat\beta = 0.478 \approx 0.5$ \textcolor{ctigreen}{\checkmark} sparse competition confirmed
\end{frame}

% ===================== SLIDE 10 =====================
\begin{frame}{The Gumbel Race: Steps 6--7 --- The Final Law}
\textbf{Gap is linear in $\kappa$:}
$\Delta_\mathrm{nearest}/s \approx a_1 \kappa_\mathrm{nearest} + a_0$

\vspace{0.5cm}
\textbf{Combine:}
\begin{block}{The CTI Universal Law}
\[
\logit(q_\mathrm{norm}) = \underbrace{\alpha}_{\text{slope}} \kappa_\mathrm{nearest} - \underbrace{\beta}_{\text{competition}} \log(K{-}1) + \underbrace{C}_{\text{offset}}
\]
\end{block}

\vspace{0.3cm}
\begin{columns}[T]
\begin{column}{0.45\textwidth}
\textbf{Derived (theorem):}
\begin{itemize}
\item Logistic in $\kappa$
\item Logarithmic in $K$
\end{itemize}
\end{column}
\begin{column}{0.45\textwidth}
\textbf{Estimated (data):}
\begin{itemize}
\item $\alpha = 1.477$ (NLP)
\item $\beta = 0.478$
\item $C$ per dataset
\end{itemize}
\end{column}
\end{columns}
\end{frame}

% ===================== SLIDE 11 =====================
\begin{frame}{Neural Collapse: Why It Works Across Architectures}
\begin{block}{Neural Collapse (Papyan et al., PNAS 2020)}
At end of training, representations converge to:
\begin{enumerate}
\item \textbf{NC1}: Within-class variability collapses
\item \textbf{NC2}: Centroids form equiangular tight frame (regular simplex)
\item \textbf{NC3}: Classifier aligns with centroids
\item \textbf{NC4}: Predictions $\to$ nearest-centroid
\end{enumerate}
\end{block}

\vspace{0.3cm}
\begin{center}
\large
$\kappa_\mathrm{nearest}$ measures \textbf{proximity to Neural Collapse}.\\[0.2cm]
Different architectures $\to$ similar NC geometry $\to$ same law.\\[0.2cm]
Architecture determines the \textbf{path}; law governs the \textbf{endpoint}.
\end{center}
\end{frame}

% ===================== SLIDE 12 =====================
\begin{frame}{The Equicorrelation $\rho$: Tightest Invariant}
\begin{columns}[T]
\begin{column}{0.52\textwidth}
$\rho$ = avg whitened cosine similarity between centroid-difference vectors.

\vspace{0.3cm}
\begin{tabular}{lc}
\toprule
\textbf{Modality} & $\boldsymbol{\rho}$ \\
\midrule
NLP decoders & 0.463 \\
Audio speech & 0.455--0.464 \\
Vision (ViT, CNN) & 0.458--0.467 \\
Mouse V1 cortex & 0.466 \\
NLP encoders & 0.446--0.456 \\
\midrule
\textbf{All 6 modalities} & \textbf{0.462} \\
\bottomrule
\end{tabular}
\end{column}
\begin{column}{0.45\textwidth}
\begin{exampleblock}{Key Facts}
\begin{itemize}
\item CV = 1.0\% (tightest invariant!)
\item Simplex prediction: $\rho = 0.5$
\item Observed: 8\% below perfect NC
\item \textbf{Substrate-independent}
\end{itemize}
\end{exampleblock}

\vspace{0.2cm}
\begin{block}{$\alpha$ prediction}
$\alpha_\mathrm{pred} = \sqrt{4/\pi}/\sqrt{1-\rho} = 1.534$\\
Observed: 1.477 (+3.8\%)\\
\textbf{Zero free parameters!}
\end{block}
\end{column}
\end{columns}
\end{frame}

% ===================== SLIDE 13 =====================
\begin{frame}{$\alpha$ by Signal Type}
\begin{center}
\begin{tabular}{lcc}
\toprule
\textbf{Signal Type} & $\boldsymbol{\alpha}$ & \textbf{Modalities} \\
\midrule
Discrete tokens & $\sim$1.48 & NLP decoders (CV = 2.3\%) \\
Continuous signals & $\sim$3.3--5.0 & ViT (4.5), CNN (4.0), Audio (4.7) \\
Encoders & $\sim$7.7 & BERT, DeBERTa, ELECTRA \\
\bottomrule
\end{tabular}
\end{center}

\vspace{0.5cm}
\begin{itemize}
\item Within NLP decoders: $\alpha$ is a \textbf{true constant} (CV consistent with estimation noise)
\item Across families: varies by signal type (NOT a failure --- level-2 universality)
\item All continuous-signal modalities converge at $\alpha \approx 4$--5
\item \textbf{Open problem}: why the gap? Same $\rho$, different $\alpha$ $\to$ $\Sigma_W$ structure
\end{itemize}
\end{frame}

% ===================== SLIDE 14 =====================
\begin{frame}{Cross-Modal Validation}
\begin{center}
\small
\begin{tabular}{lcccl}
\toprule
\textbf{Modality} & $K$ & $\alpha$ & Result & Status \\
\midrule
NLP (12 arch.) & 4--28 & 1.477 & $R^2 = 0.955$ & \textcolor{ctigreen}{\textbf{PASS}} \\
ViT-Large & 10 & 4.5 & $R^2 = 0.964$ & \textcolor{ctigreen}{\textbf{PASS}} \\
ResNet50 CNN & 100 & 4.4 & $r = 0.749$ & \textcolor{ctigreen}{\textbf{PASS}} \\
Audio (7 models) & 36 & 4.669 & $r = 0.898$ & \textcolor{ctigreen}{\textbf{PASS}} \\
Mouse V1 (32 mice) & 118 & --- & 30/32 & \textcolor{ctigreen}{\textbf{PASS}} \\
\midrule
Protein LMs & 7 & $-1.17$ & $r = -0.15$ & \textcolor{ctired}{\textbf{FAIL}} \\
Human fMRI & 12 & --- & $r = 0.12$ & \textcolor{ctired}{\textbf{FAIL}} \\
\bottomrule
\end{tabular}
\end{center}

\vspace{0.3cm}
\textbf{Form universality}: NLP, vision, audio, biological cortex. Same equation.\\
\textbf{Honest scope}: protein LMs and fMRI fail (documented in paper).
\end{frame}

% ===================== SLIDE 15 =====================
\begin{frame}{The Generation Law}
\begin{block}{Key Insight}
Next-token prediction = $V$-way classification. Same Gumbel race.\\
$\kappa$ from unembedding matrix $W_U$ (zero forward passes!) predicts perplexity.
\end{block}

\vspace{0.3cm}
\begin{center}
\begin{tabular}{lcc}
\toprule
\textbf{Test} & $r(\kappa, \log\mathrm{CE})$ & \textbf{Note} \\
\midrule
All 22 models & $-0.55$ ($p = 0.008$) & 4 arch types \\
Excl.\ LFM outlier (21) & $-0.70$ ($p < 0.001$) & \\
Fixed-$V$ Pythia+Mamba (10) & $-0.837$ ($p = 0.003$) & Arch-independent \\
\bottomrule
\end{tabular}
\end{center}

\vspace{0.3cm}
\textbf{Vocabulary drops out}: $\beta_\mathrm{gen} \approx 0$ because $K_\mathrm{eff} \approx 2$--$3$ (not $V$).\\
Transformers and SSMs lie on the \textbf{same line} ($F$-test $p = 0.147$).
\end{frame}

% ===================== SLIDE 16 =====================
\begin{frame}{Causal Evidence: 5 Tiers}
\begin{enumerate}
\item \textbf{Frozen do-interventions}: Move centroid $\to$ nearest $r = 0.899$; farthest $r = 0.000$
\item \textbf{Orthogonal factorial}: Arm A $r = 0.899$; Control $r = 0.000$
\item \textbf{Cross-modal dose-response}: 6 unseen architectures, 6/6 PASS
\item \textbf{Confusion matrix prediction (pre-registered, zero refit):}
\end{enumerate}

\begin{center}
\begin{tabular}{cccc}
\toprule
$\delta$ & $r$ & Sign accuracy & $p$-value \\
\midrule
1.0 & 0.842 & 92.9\% & $10^{-50}$ \\
2.0 & 0.808 & \textbf{100\%} & $10^{-43}$ \\
3.0 & 0.776 & \textbf{100\%} & $10^{-38}$ \\
\bottomrule
\end{tabular}
\end{center}

\begin{enumerate}
\setcounter{enumi}{4}
\item \textbf{Top-$m$ sweep}: $m{=}2$ recovers LOAO $\alpha$ within 2.3\%
\end{enumerate}
\end{frame}

% ===================== SLIDE 17 =====================
\begin{frame}{Three-Level Universality}
\begin{block}{Level 1: Form (strongest)}
Logit-linear structure is EVT-derived. Confirmed across NLP, vision, audio, biology.\\
\textit{Analogy: $PV = nRT$ holds for all ideal gases.}
\end{block}

\begin{block}{Level 2: Constant within signal type}
$\alpha$ stable within families (NLP decoders: CV = 2.3\%). Varies across signal types.\\
\textit{Analogy: van der Waals $a, b$ differ by substance.}
\end{block}

\begin{alertblock}{Level 3: Intercept is task-specific}
$C$ encodes dataset difficulty. One calibration measurement required per task.\\
\textit{Analogy: ground-state energy differs by system.}
\end{alertblock}
\end{frame}

% ===================== SLIDE 18 =====================
\begin{frame}{Open Directions}
\textbf{Near-term:}
\begin{itemize}
\item Expanded generation law (25+ models)
\item External replication
\item COLM 2026 submission (Mar 26/31)
\end{itemize}

\textbf{Open theory:}
\begin{itemize}
\item \textbf{Renormalization}: Why $\alpha_\mathrm{NLP} \neq \alpha_\mathrm{ViT}$ with same $\rho$?
\item Non-asymptotic bounds (Berry--Esseen for Gumbel)
\item Multilabel extension
\end{itemize}

\textbf{Long-term:}
\begin{itemize}
\item Training objective: maximize $\kappa_\mathrm{nearest}$
\item Connection to scaling laws: $q(C) = \sigma(\alpha C^\gamma + C_0)$
\item Universal constant from first principles
\end{itemize}
\end{frame}

% ===================== SLIDE 19 =====================
\begin{frame}{The Complete Derivation in 7 Steps}
\small
\begin{enumerate}
\item \textbf{Define} $\kappa_\mathrm{nearest}(k) = \min_{j\neq k} \|\boldsymbol{\mu}_j - \boldsymbol{\mu}_k\| / (\sigma_W\sqrt{d})$
\item \textbf{1-NN race}: correct $\iff$ $M_+ < \min_j M_j$
\item \textbf{EVT}: $M_+, M_j \to \mathrm{Gumbel}(\mu, s)$ with shared scale
\item \textbf{Gumbel race} (exact): $P = [1 + \sum \exp(-\Delta/s)]^{-1}$
\item \textbf{Sparse hierarchy}: $\sum \approx (K{-}1)^{0.5} \exp(-\Delta_\mathrm{near}/s)$
\item \textbf{Linearize}: $\Delta_\mathrm{near}/s \approx a_1\kappa + a_0$
\item \textbf{Combine}:
\end{enumerate}

\begin{center}
\LARGE\color{ctiblue}
$\logit(q) = \alpha\,\kappa - \beta\log(K{-}1) + C$
\end{center}

\normalsize
\textbf{Derived}: the functional form (logistic, logarithmic).\\
\textbf{Estimated}: the constants ($\alpha, \beta, C$).\\
\textbf{Conditional theorem}: IF assumptions hold, THEN form is necessary.
\end{frame}

% ===================== SLIDE 20 =====================
\begin{frame}{}
\begin{center}
\Huge\color{ctiblue}
One equation.\\[0.3cm]
Five modalities.\\[0.3cm]
First principles.\\[1cm]

\Large\color{black}
$\logit(q) = \alpha\,\kappa_\mathrm{nearest} - \beta\log(K{-}1) + C$
\end{center}
\end{frame}

\end{document}
